A model that answers natural-language questions is easy to build and
hard to evaluate. But there is one class of task where the evaluation is
not a judgment call at all: \textbf{write a program.} We write a
description, the model writes code, and a computer can tell us, with
certainty, whether that code does what was asked. There are no opinions
about the answer, no ambiguous grading rubric, no ``well, it's sort of
correct.'' Either the program returns the right thing for every test, or
it does not.

That one observation is why this whole book rests on the
\texttt{natural-language\ →\ Python\ program} task. It gives us
something few machine-learning setups have: \emph{objective, executable,
reproducible} evaluation. But the task alone is not enough. Before we
can run any experiment, we need a world the model can live in --- a
world that is safe, that is \emph{deterministic}, and that is small
enough that we can reason about everything that happens inside it.

\section{Why we need a world at all}\label{why-we-need-a-world-at-all}

Imagine hiring someone to write code, but not fully trusting them. We
would not hand them the keys to the server. We would give them a sandbox
--- a walled-off room where they can work, but where the dangerous tools
are simply not present.

That is exactly what we do with the model. It writes Python, and we need
to run that Python safely and deterministically, while deciding
precisely what it is \emph{allowed} to do. The model could, in
principle, write a program that reads files, connects to a network, or
deletes things. It has no reason to --- but we are not relying on its
good intentions. We are building a world where those actions are
\emph{impossible}, not just discouraged.

The deeper reason is that a controlled world makes the \emph{experiment}
clean. If the model's program could do anything --- read a file, check
the time, generate a random number --- then two runs of the same program
could give different answers, and we could never attribute a failure to
``the model didn't learn this'' versus ``the program did something
nondeterministic.'' We want the world to be so constrained that
\textbf{the only source of variability is the model itself.}

\section{The principle: whitelist, not
blocklist}\label{the-principle-whitelist-not-blocklist}

The naive way to constrain a program is a blacklist: deny
\texttt{import}, deny \texttt{eval}, deny \texttt{open}. This is the
instinct most people reach for first. But a blacklist is fundamentally
fragile --- there is always some construct we forgot. Python has an
enormous surface area, and attackers (or, in our case, just curious
models) find the gaps. Denying a hundred things is not safety; it is a
constant game of catch-up.

The robust way is the opposite: \textbf{an allowlist.} We parse the
program into its abstract syntax tree --- the structured representation
of the code, before it runs --- and walk through every node. If any node
type is not \emph{explicitly permitted}, we \textbf{reject the entire
program.} There is no ``mostly allowed.'' If the model emits even one
disallowed construct, the whole thing is refused before it ever
executes.

The difference is between a list of banned words and a grammar. A
banned-word list can be evaded by creative phrasing. A grammar says:
these are the only forms that exist. Anything outside the grammar is not
merely suspect; it is \emph{unparseable in our world.} That is the
philosophy of the sandbox.

This is the single most important design decision in the whole
apparatus, and it deserves to be stated plainly: \textbf{the sandbox
does not try to catch bad programs. It states upfront what is allowed,
and everything else is refused.}

\section{The capability spec}\label{the-capability-spec}

For the pilot, we whitelist a deliberately small subset of Python ---
small enough that we can hold it in our heads, large enough to express
real programming:

\begin{itemize}
\tightlist
\item
  literals (integers, floats), variables, assignment, \texttt{return}
\item
  arithmetic: \texttt{+\ -\ *\ /\ //\ \%\ **}
\item
  comparisons, boolean operators
\item
  \texttt{if\ /\ elif\ /\ else}, \texttt{for}, \texttt{while},
  \texttt{break}, \texttt{continue}, \texttt{pass}
\item
  lists, indexing, slicing
\item
  a small set of built-ins:
  \texttt{abs,\ min,\ max,\ sum,\ len,\ range,\ int,\ float,\ round}
\item
  simple \texttt{def} functions
\end{itemize}

Everything else is rejected: imports, file I/O, network, subprocess,
\texttt{eval}/\texttt{exec}, \texttt{class}, \texttt{try/except},
arbitrary attribute access, comprehensions, lambdas. We do not include
these because they are dangerous; we exclude them because they are not
needed for the task we are testing, and because each one we add makes
the world bigger and harder to reason about. A smaller world is a
\emph{better} world for this book, because every construct that is
allowed is one we must understand the model's relationship to.

We also version this spec (\texttt{python\_capability\_v0\_1}) so that
later, when we give the model \emph{more} Python (Chapter 18), we can
point to exactly what grew. Versioning the world is what lets us claim,
later, that a change in behavior was caused by a change in the world,
and not by something else that drifted.

\section{Resource and determinism
policy}\label{resource-and-determinism-policy}

Beyond what is allowed syntactically, we also control \emph{how} the
program may run. This is the second half of making the world
trustworthy:

\begin{itemize}
\tightlist
\item
  maximum AST depth and node count (a program can only be so complex),
\item
  maximum nesting,
\item
  a hard execution timeout (a program that loops forever is a failed
  program),
\item
  recursion protection,
\item
  float-comparison tolerance (because \texttt{0.1\ +\ 0.2} is not
  exactly \texttt{0.3}),
\item
  \textbf{no randomness, no time, no environment access.}
\end{itemize}

That last point is the quiet heart of the whole design. A program with
the same input must always produce the same output. If it does not, the
experiment is contaminated. We cannot tell whether a failure means ``the
model didn't learn this'' or ``the program rolled the dice.'' By making
the world deterministic, we eliminate that ambiguity entirely. The
book's every later claim rests on this.

\section{The pipeline}\label{the-pipeline}

\begin{verbatim}
source -> parse -> AST -> whitelist validation -> compile
        -> restricted execution -> semantic comparison
\end{verbatim}

Validation happens on the tree \emph{before} execution, so disallowed
code never runs. Execution uses a stripped set of globals that exposes
only the allowed built-ins --- there is no \texttt{\_\_import\_\_}, no
\texttt{open}, no \texttt{eval}. The program can only see the tools we
chose to give it.

\section{An honest caveat}\label{an-honest-caveat}

It is important to be clear about what this sandbox is \emph{not}. It is
a \textbf{teaching and evaluation} sandbox --- deliberately strict about
the capability subset, but not a claim of production-grade adversarial
isolation. We are not trying to secure a server against a hostile
attacker; we are trying to run generated code and measure whether it
does what we asked. An attacker with time and motivation could likely
find ways around a small AST allowlist. We state this limitation plainly
rather than overstating what the sandbox guarantees, because the book's
credibility depends on honest boundaries, not on impressive-sounding
claims.

\section{Testing the world}\label{testing-the-world}

The sandbox is the foundation of trust in every later result. If the
sandbox leaked --- allowed something disallowed, or failed to be
deterministic --- every score in the book would be suspect. So the very
first thing we do is not run any experiment at all. We \textbf{test the
world itself.}

We build a conformance suite of 31 cases: programs that use each allowed
construct and must pass, plus programs that use each disallowed
construct (\texttt{import}, \texttt{eval}, \texttt{class}, \texttt{try},
comprehensions, lambdas, arbitrary attribute access) and must all be
rejected. A sandbox that only permits things is easy; a sandbox that
also \emph{correctly refuses} everything outside its grammar is the real
achievement --- and that is what we verify before we trust a single
number it produces.

This is the first instance of a discipline the book will return to again
and again (and which we will see was worth it, painfully, in Chapter
10): \textbf{before we use a tool to measure something, we verify that
the tool measures what we think it does.} Here, the tool is the world,
and we refuse to run a single experiment until it passes its own
examination.

\section{Two gates, not one}\label{two-gates-not-one}

Everything above describes what the model is \emph{allowed to produce}.
But there is a second, quieter question: what is the model \emph{ever
asked} to do? These are not the same thing, and the difference is worth
making sharp, because it is the difference between ``the model cannot
escape'' and ``the model was never put in the position to.''

A language model has no notion of refusal. It is not a system that
decides whether a request is in scope; it is a next-token predictor.
Hand it a request outside its designed job and it does not politely
decline --- it simply keeps predicting tokens and produces whatever the
weights trend toward, often a confident near-program or garbage. So ``it
only does its one job'' can never be a property of the model's own
intent. It has to be a property of the wiring around it.

That gives us \textbf{two independent gates}, and this chapter (so far)
has really only built one of them:

\textbf{Gate A --- the output gate (the sandbox).} Already described
above: the AST allowlist rejects any generated program that uses a
construct outside the whitelist, before it runs. This guarantees that
even if the model \emph{tries} to emit \texttt{import\ os} or
\texttt{eval()}, the program is refused. Gate A is airtight on what the
model \emph{produces}.

\textbf{Gate B --- the input gate (the request validator).} This is the
gate we build for the other half. It sits \emph{in front of} the model
and decides, from the incoming natural-language request alone, whether
that request is one the model was designed for. If not, the request is
declined before the model is ever consulted.

Gate B works on the same principle as the sandbox: a \textbf{recognizer,
not a rejector of the exceptions}. It holds the set of designed
template-family shapes (the exact phrase skeletons the teacher renders)
and the designed operation vocabulary. A request is in scope only if it
fits one of those shapes \emph{and} names a known operation ---
otherwise it is refused. Just as the sandbox says ``these are the only
forms of program that exist,'' Gate B says ``these are the only forms of
request that exist.''

Why do we need both? Because each covers a failure the other cannot:

\begin{itemize}
\tightlist
\item
  \textbf{Gate A alone} would let the model be \emph{asked} for anything
  --- including a task it was never trained for --- and only then catch
  the malformed result. The model would churn, and our whole experiment
  would trust that an out-of-scope query produced a rejected program
  rather than a plausible-looking-but-wrong one.
\item
  \textbf{Gate B alone} would let the model \emph{produce} anything
  within an in-scope request, including a program that escapes its
  world, because nothing would constrain the output.
\end{itemize}

Only together do they close the loop: Gate B ensures the model is never
\emph{asked} for something outside its job, and Gate A ensures it can
never \emph{produce} something outside its world. That is the only way
to make ``this model does exactly one thing'' a true statement, rather
than an optimistic description.

There is a deep parallel here with the book's central thesis. The model
is the instrument; the experiment is the subject. In the same way, the
\emph{boundaries} are never the model --- they are always the apparatus:
the sandbox on the way out, the validator on the way in. A model that
``only does its one job'' is confined by its wiring, not by itself. That
is worth stating plainly, because it is the mistake people make when
they imagine a well-behaved model as one that simply chooses to behave.
